\chapter{Introdução}


\section{Motivação e contexto}

Os materiais compósitos são definidos pela combinação de dois ou mais constituintes macroscópicos distintos, concebida com o objetivo de obter propriedades físicas e mecânicas superiores às dos materiais homogêneos tradicionais \cite{RMJones2018}. A premissa fundamental desta classe de
materiais reside no fato de que a resposta mecânica do sistema resultante não pode ser prevista a partir de nenhum de seus constituintes isolados, pois emerge precisamente da interação entre as fases que o compõem \cite{PartonReport}. Tal característica confere aos compósitos uma capacidade singular de ajuste de propriedades  (rigidez, resistência específica, tenacidade, condutividade térmica, etc.) tornando-os indispensáveis em setores que impõem restrições simultâneas e conflitantes de desempenho e massa.

Na engenharia atual, a utilização dessas estruturas é fundamental em domínios que exigem alta resistência específica aliada à redução de peso, como as indústrias aeroespacial e naval. Neste último segmento, duas configurações estruturais concentram especial relevância tecnológica: os vasos de pressão de múltiplas camadas reforçados por fibra (\textit{Composite Overwrapped Pressure Vessels}, COPVs) \cite[p.~1]{Zhang2022} e os risers flexíveis utilizados em sistemas de produção de petróleo e gás em águas profundas \cite{Liu2026}. Aplicações críticas incluem o uso de painéis laminados em aeronaves comerciais (como no caso do Boeing 777, cuja estrutura é composta majoritariamente por compósitos) e vasos de alta pressão \cite[p.~2]{Zhang2022}. Ambas as estruturas operam sob condições severas e combinadas de carregamento (pressão interna elevada e solicitações cíclicas) o que impõe ao engenheiro a necessidade de modelos matemáticos capazes de reproduzir com fidelidade a resposta mecânica local e global do componente.

O comportamento desses materiais é governado pela interação entre suas
fases constituintes em múltiplas escalas, de modo que uma análise fidedigna deve integrar informações que se estendem desde a microestrutura  (onde se definem as propriedades de cada fase individualmente)  até o nível macroscópico, no qual o componente estrutural como um todo é dimensionado \cite{PartonReport}. Essa necessidade de trânsito entre escalas distintas configura o chamado \textbf{problema multiescala}, cujo tratamento rigoroso constitui o núcleo temático do presente trabalho.

O Método dos Elementos Finitos (MEF) constitui, atualmente, a abordagem numérica predominante na análise de estruturas compósitas. Contudo, sua aplicação direta a meios com propriedades rapidamente oscilantes  (característicos de compósitos microperiódicos) impõe uma exigência crítica: a malha de discretização deve ser suficientemente refinada para capturar as variações de propriedade em escala de $l$, o que resulta em sistemas lineares de dimensão proibitiva quando $\epsilon$ é muito pequeno \cite{Guinovart2024}.

Como alternativa matematicamente rigorosa a esta limitação, o Método de
Homogeneização Assintótica (MHA) propõe substituir o meio heterogêneo
por um material homogêneo equivalente, cujas propriedades efetivas
encapsulam os efeitos da microestrutura \cite{bakhvalov1989}.
% TRECHO ANTIGO
% Aplicações críticas incluem o uso de painéis laminados em aeronaves comerciais  como no caso do Boeing 787 Dreamliner, cuja estrutura é composta majoritariamente por compósitos  e vasos de pressão que operam sob condições severas de carregamento [Lei25, Per24]. O comportamento desses materiais é governado pela interação complexa entre suas fases, exigindo uma análise que integre dados desde a microestrutura até o nível macroscópico para prever falhas e otimizar o desempenho [Bak89, Par91].
% Dentre os tipos de compósitos, destacam-se os chamados micro-heterogêneos, cujas propriedades físicas variam de forma rápida na escala microscópica, enquanto manifestam uma aparência de homogeneidade na escala macroscópica [Lei17, Pan05]. Tais materiais podem ser encontrados tanto na natureza, como em estruturas ósseas e madeiras, quanto em meios fabricados pelo homem para aplicações de alta performance [Lei17].

\section{Problema multiescala e a homogeneização}
A modelagem de componentes micro-heterogêneos fundamenta-se no \textbf{princípio da separação de escalas} e na \textbf{hipótese do contínuo}. O princípio de separação de escalas supõe que exista uma clara distinção entre as escalas microscópica e macroscópica, sendo caracterizado por um parâmetro geométrico adimensional $\epsilon$, definido por \cite[p.~17]{bakhvalov1989} como a razão entre os comprimentos característicos da microescala $l$ (associado ao período da microestrutura) e da macroescala $L$ (associado à dimensão da estrutura como um todo). Ou seja:

\[
  \epsilon = \frac{l}{L}
\]
\begin{figure}[htbp]
  \centering
  \includegraphics[width=0.8\textwidth]{../images/separacao_de_escalas.png}
  \caption{Ilustração do processo de separação de escalas e parâmetro $\epsilon$. Fonte: \cite{BrunaReport}.}
  \label{fig:SeparacaoEscalas}
\end{figure}

De forma similar, a hipótese do contínuo estabelece que, mesmo havendo uma microestrutura bem definida em escala $l$, o material pode ser tratado como um contínuo em ambas as escalas, desde que a região de interesse seja suficientemente grande em comparação com o período da microestrutura. Pela \autoref{fig:SeparacaoEscalas}, percebe-se a relação do parâmetro $\epsilon$ com as escalas micro e macro, assim como os conceitos de separação de escalas e a hipótese do contínuo. A validade das abordagens multiescala repousa sobre a desigualdade $0 < \epsilon \ll 1$. Quando $\epsilon$ é suficientemente pequeno, torna-se possível desacoplar os problemas nas duas escalas e tratá-los de forma sequencial, reduzindo drasticamente a complexidade computacional do problema original. Para condução da análise, emprega-se uma \textit{Solução Assintótica Formal} (SAF) em série de potências do parâmetro $\epsilon$, a saber: \marginpar{\tiny Essa parte talvez fosse melhor introduzida na revisão bibliográfica}

\begin{equation}
  \mathbf{u}^{(\infty)}(\mathbf{x}, \epsilon) = \mathbf{u}^0(\mathbf{x}) + \epsilon \mathbf{u}^1(\mathbf{x}, \mathbf{y}) + \epsilon^2 \mathbf{u}^2(\mathbf{x}, \mathbf{y}) + \dots
  \label{eq:saf_expansao}
\end{equation}


onde $\mathbf{x}$ denota as coordenadas macroscópicas e $\mathbf{y} = \mathbf{x} slash \epsilon$ as coordenadas microscópicas. Estas últimas são responsáveis por descrever as oscilações rápidas da solução dentro de cada período da célula unitária \cite{bakhvalov1989}. A substituição desta expansão nas equações de equilíbrio e a coleta de termos por potências de $\epsilon$ permitem decompor o problema original em uma hierarquia de subproblemas: um \textit{problema local}, definido sobre a célula unitária e responsável por calcular os \textbf{coeficientes efetivos} do material homogeneizado, e um \textit{problema macroscópico}, formulado sobre o domínio global e governado por esses coeficientes \cite[p.~17]{BrunaReport}.

Uma distinção fundamental separa as abordagens de ordem 0 e de ordem 1 da SAF. A solução de ordem 0 fornece uma aproximação satisfatória para os campos de deslocamento macroscópico, porém falha em reproduzir com precisão as oscilações locais de tensão que caracterizam a resposta em escala de microestrutura \cite{BrunaReport}. Para superar esta limitação, o presente trabalho adota a SAF de primeira ordem, que incorpora o termo $\epsilon \mathbf{u}^1$ e é capaz de reproduzir as flutuações locais da solução exata com acurácia significativamente superior, sem aumentar proporcionalmente o custo computacional do problema macroscópico \cite{BrunaReport}. \marginpar{\tiny Passar os trechos da homogeneização para a subseção motivação e contexto, acho que vai fazer mais sentido para contrasta-lo com o MEF}

\section{Objetivo geral}

O objetivo geral deste trabalho é realizar uma análise multiescala de
vasos de pressão compósitos microperiódicos,  estabelecendo um comparativo
sistemático de precisão e eficiência computacional entre o Método dos
Elementos Finitos (MEF), empregado como solução de referência, e o Método
de Homogeneização Assintótica de primeira ordem (MHA), empregado como
método principal de análise.

\section{Objetivos específicos}
Os objetivos específicos que sustentam esta meta são os seguintes:

\begin{enumerate}
  \item \textbf{Formular o problema de equilíbrio do cilindro compósito microperiódico}
  em coordenadas cilíndricas, sob hipótese de axissimetria, e reduzi-lo à
  forma escalar unidimensional que serve de ponto de partida para a
  aplicação do MHA % \cite{BrunaReport} \cite{bakhvalov1989}.
\end{enumerate}

\begin{enumerate}
  \item \textbf{Aplicar o MHA ao problema formulado}, conduzindo a dedução formal da
  expansão assintótica, dos problemas locais sobre a célula unitária e da
  equação homogeneizada que governa o deslocamento macroscópico.
% \cite{bakhvalov1989} \cite{PartonReport}
\end{enumerate}

\begin{enumerate}
  \item \textbf{Deduzir e calcular os coeficientes efetivos de elasticidade} do
  material equivalente para a configuração axissimétrica bilaminada,
  expressando-os em termos das frações volumétricas e das propriedades
  constitutivas de cada fase. % \cite{BrunaReport}
\end{enumerate}

\begin{enumerate}
  \item \textbf{Desenvolver e validar um código computacional} baseado na biblioteca
  FEniCS para a simulação numérica do problema de elasticidade
  axissimétrica via MEF, permitindo a obtenção de soluções de referência
  para a comparação com os resultados do MHA.
\end{enumerate}

\begin{enumerate}
  \item \textbf{Comparar quantitativamente a precisão e a eficiência computacional}
  entre a solução homogeneizada de segunda ordem e a solução numérica
  detalhada, avaliando o impacto da inclusão dos termos de ordem superior
  na SAF sobre a qualidade da aproximação dos campos de tensão e
  deslocamento % \cite{BrunaReport}.
\end{enumerate}


\section{Organização e limitações do Trabalho}

O presente trabalho está organizado em seis capítulos. O \textbf{Capítulo 2}
apresenta a revisão bibliográfica, cobrindo os fundamentos dos meios
micro-heterogêneos, as equações de equilíbrio mecânico em coordenadas
cilíndricas sob hipótese de axissimetria, as relações constitutivas da
elasticidade linear isotrópica e os princípios teóricos do MHA e do MEF,
incluindo uma descrição da plataforma computacional FEniCS. O \textbf{Capítulo 3}
detalha os materiais e métodos, apresentando a formulação completa do
problema de valor de contorno, a dedução da SAF, o procedimento de
homogeneização e o cálculo analítico dos coeficientes efetivos. O
\textbf{Capítulo 4} expõe os resultados obtidos e a discussão comparativa entre
as abordagens numéricas e analíticas. O \textbf{Capítulo 5} consolida as
conclusões do trabalho e aponta perspectivas para investigações futuras.
As referências bibliográficas são listadas ao final do documento.

Neste trabalho não será abordado ...

\marginpar{\tiny Lembrete: Falar sobre o escopo do trabalho. O que não será abordado}
